\SourceLeaf{367}%
und ein kurzer Auszug daraus wurde schon damals in unserer Zeitschrift mitgeteilt (1805, Stück 124). Wir bemerken hier also nur, dass es sich um die Stelle handelt, an der Sokrates durch ein Beispiel aus der Geometrie anschaulich machen will, wie man sich auf die Lösung einer Aufgabe vorbereitet, indem man vorher gewisse genauere Bestimmungen annimmt. Die geometrische Aufgabe, die Sokrates dafür wählt, ist die Frage, ob sich ein gegebenes Dreieck einem gegebenen Kreis einbeschreiben lässt. Die Worte aber, mit denen er zunächst gewisse Einschränkungen für die Art des Dreiecks festlegen will, haben den Auslegern viel Mühe bereitet. Herr \textsc{Mollweide} zeigt mit großem gelehrtem Scharfsinn, dass die bezeichnete Eigenschaft darin besteht, dass sich das Dreieck in zwei andere zerlegen lässt, die dem ganzen Dreieck ähnlich sind. Das ist freilich im Grunde nur eine gekünstelte Umschreibung des \textit{rechtwinkligen} Dreiecks. Herr M. beweist, dass diese Eigenschaft nur dem rechtwinkligen Dreieck zukommen kann, auf eine viel kunstvollere und umständlichere Weise, als hier nötig gewesen wäre. Das folgt nämlich unmittelbar aus der Gleichheit der drei Winkel \(ABC,ADB,BDC\) (S.~46).

Auch die dritte Abhandlung war unserer Sozietät schon früher als Manuskript vorgelegt worden, und in unseren \textit{gelehrten Anzeigen} war darüber berichtet worden (1807, Stück 74). Sie erscheint hier mit erheblichen Ergänzungen. Darin werden zwei von \textsc{Columella} gelehrte Näherungsverfahren zur Berechnung des Flächeninhalts eines gleichseitigen Dreiecks und eines Kreissegments erläutert. Auf S.~71 findet sich eine kleine Übereilung. Dort wird behauptet, kein anderer Bruch, dessen Zähler und Nenner unter 100 liegen, komme dem wahren Verhältnis des gleichseitigen Dreiecks zum Quadrat über derselben Seite so nahe wie \(\frac{13}{30}\). Tatsächlich sind die beiden Brüche \(\frac{29}{67}\) und \(\frac{42}{97}\) genauer.

Der Brief an den verdienten Professor \textsc{Schneider} in Breslau enthält einige Anmerkungen zu den Auszügen, die \textsc{Hase} in \textsc{Bredow}s \textit{Epistolae Parisienses} aus den freilich sehr unbedeutenden mathematischen Schriften des \textsc{Vitruvius Rufus} und des \textsc{Epaphroditus} mitgeteilt hat.
